\vspace{-20pt}
\section{Introduction}
\label{sec:intro}

Effective communication is a key ability for collaboration.
Indeed, intelligent agents (humans or artificial) in real-world scenarios can
significantly benefit from exchanging information that enables them to coordinate,
strategize, and utilize their combined sensory experiences to act in the physical world.
The ability to communicate has wide-ranging applications for artificial agents --
from multi-player gameplay in simulated (\eg DoTA, StarCraft) or
physical worlds (\eg robot soccer), to self-driving car networks communicating
with each other to achieve safe and swift transport, to teams of robots on
search-and-rescue missions deployed in hostile, fast-evolving environments.

A salient property of human communication is the ability to hold \emph{targeted} interactions.
Rather than the `one-size-fits-all' approach of broadcasting messages to all participating agents,
as has been previously explored~\citep{sukhbaatar_nips16,foerster_nips16,singh_iclr19},
it can be useful to direct certain messages to specific recipients.
This enables a more flexible collaboration strategy in complex environments.
For example, within a team of search-and-rescue robots with a diverse set of roles and goals,
a message for a fire-fighter (\eg ``smoke is coming from the kitchen'') is largely
meaningless for a bomb-defuser.

\begin{table*}[t]
    \vspace{-5pt}
    \setlength\tabcolsep{6pt}
    \centering
    \begin{tabular}{@{}l c c c r@{}}
        &  \footnotesize Decentralized  & \footnotesize Targeted & \footnotesize Multi-Round & \footnotesize Reinforcement  \\[-3pt]
        & \footnotesize Execution & \footnotesize  Communication & \footnotesize Decisions & \footnotesize Learning \\
        \cmidrule(r){1-1} \cmidrule(lr){2-2} \cmidrule(lr){3-3} \cmidrule(lr){4-4} \cmidrule(lr){5-5}
        DIAL~\citep{foerster_nips16}
            & Yes & No & No & Yes (Q-Learning) \\
        CommNet~\citep{sukhbaatar_nips16}
            & Yes & No & Yes & Yes (REINFORCE) \\
        VAIN~\citep{hoshen_nips17}
            & No & Yes & Yes & No (Supervised) \\
        ATOC~\citep{jiang_nips18}
            & Yes & No & No & Yes (Actor-Critic) \\
        IC3Net~\citep{singh_iclr19}
            & Yes & No & Yes & Yes (REINFORCE) \\ \midrule
        TarMAC (this paper) & Yes & Yes & Yes & Yes (Actor-Critic) \\ \bottomrule
    \end{tabular}
    \vspace{5pt}
    \caption{Comparison with previous work on collaborative multi-agent communication with continuous vectors.}
    \vspace{-15pt}
    \label{table:related}
\end{table*}

We develop TarMAC, a Targeted Multi-Agent Communication
architecture for collaborative multi-agent deep reinforcement learning.
Our key insight in TarMAC is to allow each individual agent to \emph{actively select}
which other agents to address messages to. This targeted communication
behavior is operationalized via a simple signature-based soft attention mechanism: along with the message, the sender broadcasts
a key which encodes properties of agents the message is intended for, and is used by receivers to gauge the relevance of the message.
This communication mechanism is learned implicitly, without any attention supervision,
as a result of end-to-end training using task reward.

The inductive bias provided by soft attention in the communication architecture is sufficient to enable agents to 1) communicate agent-goal-specific messages (\eg guide fire-fighter towards fire, bomb-defuser towards
bomb, \etc), 2) be adaptive to variable team sizes (\eg the size of the local neighborhood a self-driving car can communicate with changes as it moves), and 3) be interpretable through predicted attention probabilities that allow for inspection of \emph{which} agent is communicating \emph{what} message and to \emph{whom}.

Our results however show that just using targeted communication is not enough.
Complex real-world tasks might require \emph{large populations of agents} to go through
\emph{multiple rounds of collaborative communication and reasoning}, involving
large amounts of information to be \emph{persistent in memory} and exchanged via
\emph{high-bandwidth communication channels}. To this end, our actor-critic
framework combines centralized training with decentralized execution~\citep{lowe_nips17},
thus enabling scaling to large team sizes. In this context, our inter-agent
communication architecture also supports multiple rounds of targeted interactions at
every time-step, wherein the agents' recurrent policies persist
relevant information in internal states.

While natural language, \ie a finite set of discrete tokens with pre-specified human-conventionalized meanings, may seem like an intuitive protocol for inter-agent communication -- one that enables human-interpretability of interactions -- forcing machines to communicate among themselves in discrete tokens presents additional training challenges. Since our work focuses on machine-only multi-agent teams, we allow agents to communicate via continuous vectors (rather than discrete symbols),
as has been explored in~\cite{sukhbaatar_nips16,singh_iclr19},
and agents have the flexibility to discover and optimize their communication protocol as per task requirements.

We provide extensive empirical evaluation of our approach across
a range of tasks, environments, and team sizes.
\begin{compactitem}
\item We begin by benchmarking TarMAC and its ablation without attention on a
    cooperative navigation task derived from the SHAPES environment~\citep{andreas_cvpr16} in \Secref{sec:shapes}.
    We show that agents learn intuitive attention behavior across task difficulties.
\item Next, we evaluate TarMAC on the traffic junction environment~\citep{sukhbaatar_nips16} in \Secref{sec:tj},
    and show that agents are able to adaptively focus on `active' agents in the
    case of varying team sizes.
\item We then demonstrate its efficacy in $3$D environments with a cooperative first-person point-goal navigation task in House$3$D~\citep{house3d} (\Secref{sec:h3d}).
\item Finally, in \Secref{sec:competitive}, we show that TarMAC can be easily combined with IC3Net~\cite{singh_iclr19},
    thus extending its applicability to mixed and competitive environments, and
    leading to significant improvements in performance and sample complexity.
\end{compactitem}
